Ranked by mean single-electrode macro $F_1$, Fp1, Fp2, AF4 and AF3 topped Internal's ranking, reaching $F_1$ values of 83.73\%, 83.37\%, 75.16\% and 67.59\%, respectively, the same four frontal electrodes identified by the earlier restricted analysis and in the same order except that Fp1 now leads Fp2 by 0.36 percentage points rather than trailing it (Figure~\ref{fig:exp1_single_channel_stats}; Figure~\ref{fig:exp1_single_channel}). C3 followed next in the ranking at $F_1=54.05\%$, ahead of the remaining two canonical frontal electrodes, F4 and F3, which trailed further behind at $39.98\%$ and $39.85\%$. Performance then fell sharply across the remaining 17 electrodes. The weakest was C4 at $F_1=1.51\%$, with every parietal and occipital electrode below $3.2\%$.

On Cao2018, Fp1, Fp2, F3 and F4 topped the single-electrode ranking, the same as Internal's canonical frontal set, reaching $F_1$ values of 77.65\%, 75.59\%, 49.30\% and 49.03\%, respectively, matching both the identity and the order of the full-montage ranking (Figure~\ref{fig:exp1_single_channel_stats}). Performance then dropped substantially at the next electrode in the ranking, FC3, to $F_1=38.44\%$, followed by the remaining central electrodes FC4, FT8 and FT7, reaching $F_1$ values of 36.72\%, 34.99\% and 34.22\%, respectively. The central region therefore occupied the entire second tier of this ranking, clearly below the four frontal electrodes and still ahead of the weaker central pair, C3 and C4, at $25.43\%$ and $25.10\%$. Beyond the central electrodes, performance declined further toward parietal ($14.80$--$16.22\%$) and occipital sites, with the weakest electrode, O1, at $F_1=2.28\%$. This confirms, within the four curated regions, the same anterior-to-posterior decline reported for the full-montage and regional-subset analyses (Sections~\ref{sec:exp1_reference} and~\ref{sec:exp1_regional}).

%Comparing this complete single-electrode ranking against the full-montage per-electrode ranking (Section~\ref{sec:exp1_reference}) shows the underlying spatial pattern is essentially unchanged by removing channel context. The same four electrodes occupy the top four positions in both corpora, and Cao2018's ranking order is identical. On Internal, Fp1 and Fp2 are close enough, a 0.36-point gap under single-electrode operation versus a 0.50-point gap under the full montage, in opposite directions, that their relative order is not a meaningful reordering, consistent with neither electrode's single-versus-full difference reaching significance below. The Internal C3--C4 asymmetry, previously seen under the full montage at $61.06\%$ versus $1.85\%$ and the central subset's own rerun at $58.59\%$ versus $1.73\%$, persists when C3 and C4 operate as isolated single-electrode detectors at $54.05\%$ versus $1.51\%$, if anything widening in relative terms. Cao2018's C3/C4 pair remains close together at $25.43\%$ and $25.10\%$, confirming the asymmetry stays specific to Internal. Posterior electrodes remain the weakest condition in every corpus at every level of channel context.
